\chapter{Appendix: Sets, Functions, and Fields}

This appendix is the toolbox under the workbench. Nothing here is ``linear algebra''
proper, but everything here is the vocabulary linear algebra is spoken in. When the main
chapters say ``the set of all vectors $\vec{v}$ such that $\mat{A}\vec{v} = \vec{0}$,'' or
``the function $T$ is one-to-one,'' or ``let the scalars come from a field $\F$,'' they
are quietly assuming you are comfortable with sets, functions, and fields. Most readers
have met these ideas before in scattered places. The goal here is to gather them into one
calm, geometry-first reference you can return to whenever a definition in the main text
feels like it skipped a step.

You do not need to read this front to back. Skim it, notice what is familiar, and circle
back to whatever a later chapter leans on. But if any of the following make you uneasy --
what exactly a function ``onto'' means, why $\R$ and $\C$ get to be the scalars and the
integers $\Z$ do not, how to read a proof without panicking -- spend twenty unhurried
minutes here. It pays for itself many times over.

\chapterline

\section{Sets: the bags everything lives in}

\subsection{What a set is, and the three things you do with it}

A \emph{set} is a collection of distinct objects, called its \emph{elements} or
\emph{members}. That is the entire definition, and its plainness is the point. A set does
not care about order, does not care about repetition, and does not (by itself) come with
any operations. It is just a bag with things in it.

We write $x \in S$ to say ``$x$ is an element of $S$,'' and $x \notin S$ for ``$x$ is
not.'' Two sets are \emph{equal} exactly when they have the same elements -- nothing more
is required:
\[
  \{1,2,3\} = \{3,1,2\} = \{1,1,2,3\}.
\]
Order and duplicates are invisible to a set.

\begin{intuitionbox}
Picture a set as a loop drawn on paper with dots inside it. The loop's shape, the
position of the dots, the order you drew them -- all irrelevant. The only question a set
ever answers is: ``is this thing in, or out?'' Everything else in this section is
bookkeeping built on that one yes/no question.
\end{intuitionbox}

There are two standard ways to name a set. You can \emph{list} its elements,
\[
  S = \{2, 4, 6, 8\},
\]
or you can describe them with a \emph{rule}, read ``the set of all $x$ such that\ldots'':
\[
  S = \set{x \in \N : x \text{ is even and } 2 \le x \le 8}.
\]
The colon (some authors use a vertical bar $\mid$) means ``such that.'' The part before
it says where $x$ lives; the part after it is the membership test. This builder notation
is everywhere in linear algebra. The null space of a matrix, for instance, is
\[
  \nullsp(\mat{A}) = \set{\vec{x} \in \R^n : \mat{A}\vec{x} = \vec{0}},
\]
``the set of all vectors $\vec{x}$ in $\R^n$ that $\mat{A}$ sends to zero.'' If you can
read that line aloud and know what it is collecting, you can read most of this book.

\subsection{Numbers you already trust, named}

A handful of number sets recur so often they get their own blackboard letters. You have
seen them; here they are in one place, each nested inside the next.

\begin{center}
\begin{tabular}{ll}
\toprule
$\N = \{0, 1, 2, 3, \dots\}$ & the \emph{natural numbers} (counting numbers) \\
$\Z = \{\dots, -2, -1, 0, 1, 2, \dots\}$ & the \emph{integers} \\
$\Q = \set{p/q : p,q \in \Z,\ q \ne 0}$ & the \emph{rationals} (fractions) \\
$\R$ & the \emph{reals} (the whole number line) \\
$\C = \set{a + bi : a,b \in \R}$ & the \emph{complex} numbers \\
\bottomrule
\end{tabular}
\end{center}

Each set sits inside the next: $\N \subset \Z \subset \Q \subset \R \subset \C$
(Figure~\ref{ch38:numbernest}). The
symbol $\subset$ means \emph{subset} -- every element of the left set is also an element
of the right one. (Some authors write $\subseteq$ to stress that equality is allowed; we
use $\subset$ to mean ``is contained in,'' allowing equality, throughout.) The jump from
$\Q$ to $\R$ fills in the ``gaps'' like $\sqrt{2}$ and $\pi$ that no fraction reaches; the
jump from $\R$ to $\C$ adds a square root of $-1$. We will lean on exactly two of these,
$\R$ and $\C$, as our scalars, and Section~\ref{ch38:fields} explains the precise reason
those two are the comfortable choices.

\begin{figure}[H]
\centering
\begin{tikzpicture}[scale=1.0]
  \foreach \r in {2.6, 2.05, 1.5, 0.95, 0.45}{
    \draw[thick, structurepurple] (0,0) ellipse ({\r} and {\r*0.82});
  }
  \node[structurepurple] at (0,0.02) {\small$\N$};
  \node[structurepurple] at (0,0.6) {\small$\Z$};
  \node[structurepurple] at (0,1.1) {\small$\Q$};
  \node[structurepurple] at (0,1.62) {\small$\R$};
  \node[structurepurple] at (0,2.05) {\small$\C$};
  \node[vecred] at (1.3,-0.65) {\scriptsize $\sqrt 2,\ \pi$};
  \node[vecblue] at (1.95,-1.35) {\scriptsize $i,\ 3{+}2i$};
\end{tikzpicture}
\caption{The standard number systems nest like Russian dolls:
$\N \subset \Z \subset \Q \subset \R \subset \C$. Each layer adds something the previous
one was missing -- negatives, then fractions, then irrationals like $\sqrt 2$ and $\pi$,
then a square root of $-1$.}
\label{ch38:numbernest}
\end{figure}

A couple of special sets round out the vocabulary. The \emph{empty set} $\varnothing =
\{\,\}$ has no elements at all; it is the set-theoretic zero, and it turns up whenever a
membership test is never satisfied (a system with no solution has an empty solution set).
And for any set $S$ we may ask its \emph{cardinality} $\abs{S}$, the number of elements in
it: $\abs{\{2,4,6,8\}} = 4$, while $\abs{\R}$ is infinite.

\subsection{Combining sets}

Given two sets $A$ and $B$ living inside some common universe, three operations let you
build new sets. Each has a one-line definition and a one-picture meaning
(Figure~\ref{ch38:venn}).

\begin{itemize}
  \item \textbf{Union} $A \cup B = \set{x : x \in A \text{ or } x \in B}$ -- everything in
        either bag. (``Or'' in mathematics is always inclusive: in one, the other, or
        both.)
  \item \textbf{Intersection} $A \cap B = \set{x : x \in A \text{ and } x \in B}$ --
        everything in both bags at once.
  \item \textbf{Difference} $A \setminus B = \set{x : x \in A \text{ and } x \notin B}$ --
        what is left of $A$ after deleting anything that is also in $B$.
\end{itemize}

\begin{figure}[H]
\centering
\begin{tikzpicture}[scale=0.85]
  % Union
  \begin{scope}[xshift=0cm]
    \fill[sourceblue!25] (-0.5,0) circle (0.9);
    \fill[sourceblue!25] (0.5,0) circle (0.9);
    \draw[thick] (-0.5,0) circle (0.9);
    \draw[thick] (0.5,0) circle (0.9);
    \node at (0,-1.55) {\small $A \cup B$};
    \node at (-0.75,0) {\scriptsize$A$};
    \node at (0.75,0) {\scriptsize$B$};
  \end{scope}
  % Intersection
  \begin{scope}[xshift=4.2cm]
    \begin{scope}
      \clip (-0.5,0) circle (0.9);
      \fill[targetgreen!35] (0.5,0) circle (0.9);
    \end{scope}
    \draw[thick] (-0.5,0) circle (0.9);
    \draw[thick] (0.5,0) circle (0.9);
    \node at (0,-1.55) {\small $A \cap B$};
    \node at (-0.75,0) {\scriptsize$A$};
    \node at (0.75,0) {\scriptsize$B$};
  \end{scope}
  % Difference
  \begin{scope}[xshift=8.4cm]
    \begin{scope}
      \clip (-0.5,0) circle (0.9);
      \fill[quotientorange!35] (-0.5,0) circle (0.9);
    \end{scope}
    \begin{scope}
      \clip (0.5,0) circle (0.9);
      \fill[white] (-2,-2) rectangle (2,2);
    \end{scope}
    \draw[thick] (-0.5,0) circle (0.9);
    \draw[thick] (0.5,0) circle (0.9);
    \node at (0,-1.55) {\small $A \setminus B$};
    \node at (-0.75,0) {\scriptsize$A$};
    \node at (0.75,0) {\scriptsize$B$};
  \end{scope}
\end{tikzpicture}
\caption{The three ways to combine two sets, shaded. Union keeps everything in either
circle, intersection keeps the overlap, difference keeps the part of $A$ outside $B$.}
\label{ch38:venn}
\end{figure}

These appear in disguise all through the book. Two lines in the plane that cross meet in
the \emph{intersection} of their solution sets -- a single point. Two lines that are
parallel and distinct have empty intersection -- the system has no solution. The span of a
collection of vectors is built so that it is closed under combination, and subspaces will
turn out to be closed under intersection but \emph{not} under union, a subtlety the
chapter on vector spaces makes much of.

\subsection{Ordered pairs and the product $\R^n$}

A set forgets order, but geometry cannot afford to: the point $(1,2)$ is not the point
$(2,1)$. So we need a companion idea, the \emph{ordered pair} $(a,b)$, where position
carries meaning and $(a,b) = (c,d)$ only when $a=c$ \emph{and} $b=d$.

The \emph{Cartesian product} of two sets collects all ordered pairs:
\[
  A \times B = \set{(a,b) : a \in A,\ b \in B}.
\]
The most important instance in all of this book is
\[
  \R^2 = \R \times \R = \set{(x,y) : x,y \in \R},
\]
the plane, and more generally $\R^n$, the set of ordered lists of $n$ real numbers. That
little superscript $n$ in $\R^n$ is a Cartesian product done $n$ times. Every vector you
draw, every data point, every state of a system lives in some $\R^n$. The whole subject
is, in one sense, the study of structure-preserving maps between these product sets.

\begin{figure}[H]
\centering
\begin{tikzpicture}[scale=0.95]
  \draw[step=1, gray!35, very thin] (-0.4,-0.4) grid (3.4,2.4);
  \draw[-{Latex[length=2mm]}, thick] (-0.6,0) -- (3.7,0) node[right]{\scriptsize $x$};
  \draw[-{Latex[length=2mm]}, thick] (0,-0.6) -- (0,2.7) node[above]{\scriptsize $y$};
  % highlight the slice x=3 (a copy of R) and y=2 (a copy of R)
  \draw[sourceblue, thick] (3,-0.4) -- (3,2.4);
  \draw[targetgreen, thick] (-0.4,2) -- (3.4,2);
  \node[circle,fill=vecred,inner sep=1.6pt] at (3,2) {};
  \node[vecred, right=2pt] at (3,2) {\scriptsize $(3,2)$};
  \node[sourceblue] at (3.25,0.55) {\scriptsize $x=3$};
  \node[targetgreen] at (0.95,2.25) {\scriptsize $y=2$};
\end{tikzpicture}
\caption{The product $\R^2 = \R \times \R$ as a grid. Fixing the first coordinate at
$x=3$ (blue) traces out a vertical copy of $\R$; fixing the second at $y=2$ (green) traces
a horizontal copy. Their unique crossing is the ordered pair $(3,2)$ -- position carries
meaning, so $(3,2) \ne (2,3)$.}
\label{ch38:product}
\end{figure}

\begin{focusbox}
The four set ideas you actually need downstream: (1) builder notation
$\set{x : \text{rule}}$ to \emph{define} subspaces and solution sets; (2) subset $\subset$
to say one space sits inside another; (3) intersection $\cap$ for ``solves both
conditions''; (4) the product $\R^n$ as the home of vectors
(Figure~\ref{ch38:product}). If those four feel automatic,
the rest of this section is just polish.
\end{focusbox}

\section{Functions: machines that move points}

\subsection{Domain, codomain, rule}

A \emph{function} $f$ from a set $A$ to a set $B$, written
\[
  f : A \to B,
\]
is a rule that assigns to \emph{each} element $a \in A$ exactly \emph{one} element $f(a)
\in B$. The set $A$ is the \emph{domain} (the legal inputs), and $B$ is the \emph{codomain}
(the declared pool of possible outputs). The two words ``each'' and ``exactly one'' carry
all the weight: every input must get an answer, and no input may get two.

\begin{intuitionbox}
A function is a machine. You feed in one item from the input bag $A$; out comes one item
from the output bag $B$. The machine is deterministic (same input, same output every
time) and total (it never jams on a legal input). That is the whole contract. A matrix
$\mat{A}$ acting by $\vec{x} \mapsto \mat{A}\vec{x}$ is exactly such a machine, with
domain $\R^n$ and codomain $\R^m$.
\end{intuitionbox}

Notice the asymmetry. Every input is used exactly once, but outputs are under no such
rule: an output in $B$ may be hit by several inputs, or by none at all. Those two
loopholes -- ``hit more than once'' and ``hit not at all'' -- are precisely what the next
two definitions police.

The \emph{image} (or \emph{range}) of $f$ is the set of outputs that actually occur:
\[
  \operatorname{im}(f) = \set{f(a) : a \in A} \subset B.
\]
The image can be smaller than the codomain. The squaring function $f:\R\to\R$, $f(x)=x^2$,
has codomain $\R$ but image only $[0,\infty)$ -- negative numbers are declared possible
outputs yet never produced. For a matrix, the image is exactly the \emph{column space}:
the set of all $\vec{b}$ you can actually reach as $\mat{A}\vec{x}$.

\subsection{Injective, surjective, bijective}

These three words sound forbidding and mean something you can see in a picture
(Figure~\ref{ch38:inj-surj-bij}). They
describe how the arrows of a function land on the codomain.

\begin{figure}[H]
\centering
\begin{tikzpicture}[scale=0.92, dot/.style={circle,fill,inner sep=1.4pt}]
  % Injective not surjective
  \begin{scope}[xshift=0cm]
    \node[anchor=south,font=\footnotesize\scshape,text=spacegray!90!black] at (0,2.5) {1. Injective (One-to-One)};
    \draw[fill=codebg,draw=spacegray!30,rounded corners=4pt] (-1.8,-2.2) rectangle (1.8,2.4);

    % domain A and codomain B capsules
    \draw[fill=sourceblue!10,draw=sourceblue!50,rounded corners=6pt] (-1.4,-1.6) rectangle (-0.6,1.6);
    \draw[fill=targetgreen!10,draw=targetgreen!50,rounded corners=6pt] (0.6,-1.8) rectangle (1.4,1.8);
    \node[font=\scriptsize,text=sourceblue!90!black] at (-1.0,-1.9) {$A$};
    \node[font=\scriptsize,text=targetgreen!90!black] at (1.0,-2.0) {$B$};

    \node[dot,sourceblue] (a1) at (-1.0,1.0) {};
    \node[dot,sourceblue] (a2) at (-1.0,0.0) {};
    \node[dot,sourceblue] (a3) at (-1.0,-1.0) {};

    \node[dot,targetgreen] (b1) at (1.0,1.4) {};
    \node[dot,targetgreen] (b2) at (1.0,0.5) {};
    \node[dot,targetgreen] (b3) at (1.0,-0.4) {};
    \node[dot,targetgreen] (b4) at (1.0,-1.3) {};

    \draw[-{Latex[length=2mm]},thick,spacegray!80!black] (a1) -- (b1);
    \draw[-{Latex[length=2mm]},thick,spacegray!80!black] (a2) -- (b2);
    \draw[-{Latex[length=2mm]},thick,spacegray!80!black] (a3) -- (b3);

    \node[font=\tiny,align=center,text=spacegray!90!black] at (0,2.0) {No target hit twice\\(target missed $\Rightarrow$ not surj)};
  \end{scope}

  % Surjective not injective
  \begin{scope}[xshift=4.4cm]
    \node[anchor=south,font=\footnotesize\scshape,text=spacegray!90!black] at (0,2.5) {2. Surjective (Onto)};
    \draw[fill=codebg,draw=spacegray!30,rounded corners=4pt] (-1.8,-2.2) rectangle (1.8,2.4);

    % domain A and codomain B capsules
    \draw[fill=sourceblue!10,draw=sourceblue!50,rounded corners=6pt] (-1.4,-1.8) rectangle (-0.6,1.8);
    \draw[fill=targetgreen!10,draw=targetgreen!50,rounded corners=6pt] (0.6,-1.3) rectangle (1.4,1.3);
    \node[font=\scriptsize,text=sourceblue!90!black] at (-1.0,-2.0) {$A$};
    \node[font=\scriptsize,text=targetgreen!90!black] at (1.0,-1.6) {$B$};

    \node[dot,sourceblue] (c1) at (-1.0,1.3) {};
    \node[dot,sourceblue] (c2) at (-1.0,0.4) {};
    \node[dot,sourceblue] (c3) at (-1.0,-0.4) {};
    \node[dot,sourceblue] (c4) at (-1.0,-1.3) {};

    \node[dot,targetgreen] (d1) at (1.0,0.8) {};
    \node[dot,targetgreen] (d2) at (1.0,-0.6) {};

    \draw[-{Latex[length=2mm]},thick,spacegray!80!black] (c1) -- (d1);
    \draw[-{Latex[length=2mm]},thick,spacegray!80!black] (c2) -- (d1);
    \draw[-{Latex[length=2mm]},thick,spacegray!80!black] (c3) -- (d2);
    \draw[-{Latex[length=2mm]},thick,spacegray!80!black] (c4) -- (d2);

    \node[font=\tiny,align=center,text=spacegray!90!black] at (0,2.0) {Every target reached\\(multiple hits $\Rightarrow$ not inj)};
  \end{scope}

  % Bijective
  \begin{scope}[xshift=8.8cm]
    \node[anchor=south,font=\footnotesize\scshape,text=spacegray!90!black] at (0,2.5) {3. Bijective (Isomorphism)};
    \draw[fill=codebg,draw=spacegray!30,rounded corners=4pt] (-1.8,-2.2) rectangle (1.8,2.4);

    % domain A and codomain B capsules
    \draw[fill=sourceblue!10,draw=sourceblue!50,rounded corners=6pt] (-1.4,-1.6) rectangle (-0.6,1.6);
    \draw[fill=targetgreen!10,draw=targetgreen!50,rounded corners=6pt] (0.6,-1.6) rectangle (1.4,1.6);
    \node[font=\scriptsize,text=sourceblue!90!black] at (-1.0,-1.9) {$A$};
    \node[font=\scriptsize,text=targetgreen!90!black] at (1.0,-1.9) {$B$};

    \node[dot,sourceblue] (e1) at (-1.0,1.0) {};
    \node[dot,sourceblue] (e2) at (-1.0,0.0) {};
    \node[dot,sourceblue] (e3) at (-1.0,-1.0) {};

    \node[dot,targetgreen] (g1) at (1.0,1.0) {};
    \node[dot,targetgreen] (g2) at (1.0,0.0) {};
    \node[dot,targetgreen] (g3) at (1.0,-1.0) {};

    \draw[-{Latex[length=2mm]},thick,spacegray!80!black] (e1) -- (g1);
    \draw[-{Latex[length=2mm]},thick,spacegray!80!black] (e2) -- (g2);
    \draw[-{Latex[length=2mm]},thick,spacegray!80!black] (e3) -- (g3);

    \node[font=\tiny,align=center,text=spacegray!90!black] at (0,2.0) {Perfect 1-to-1 matching\\(invertible inverse $f^{-1}$ exists)};
  \end{scope}
\end{tikzpicture}
\caption{The three structural archetypes of a mapping. Left: Injective (no two domain elements map to the same target; $\nullsp\mat{A} = \{\vec{0}\}$). Middle: Surjective (the entire codomain is covered; $\col\mat{A} = B$). Right: Bijective (both injective and surjective -- a perfect matching, representing an invertible transformation or isomorphism).}
\label{ch38:inj-surj-bij}
\end{figure}

\begin{formalbox}
Let $f : A \to B$.
\begin{itemize}
  \item $f$ is \textbf{injective} (one-to-one) if distinct inputs give distinct outputs:
        $a_1 \ne a_2 \implies f(a_1) \ne f(a_2)$. Equivalently (and this is the form you
        use in proofs), $f(a_1) = f(a_2) \implies a_1 = a_2$. No output is hit twice.
  \item $f$ is \textbf{surjective} (onto) if every element of the codomain is hit:
        for every $b \in B$ there is some $a \in A$ with $f(a) = b$. The image fills all
        of $B$. No output is missed.
  \item $f$ is \textbf{bijective} if it is both injective and surjective: a perfect
        one-to-one pairing between $A$ and $B$, with nothing doubled and nothing left out.
\end{itemize}
\end{formalbox}

\begin{trapbox}
``Injective'' and ``surjective'' are statements about a function \emph{together with its
declared codomain}, not about the formula alone. The map $f:\R\to\R$, $f(x)=e^x$, is not
surjective ($-1$ is never an output). But the very same formula seen as
$f:\R\to(0,\infty)$ \emph{is} surjective. Change the codomain and you can change the
verdict. Always ask ``into which $B$?'' before you answer.
\end{trapbox}

\begin{example}[Testing all three on small maps]\label{ch38:smallmaps}
Take inputs and outputs in $\R$ unless noted.
\begin{itemize}
  \item $f(x) = 2x$ on $\R \to \R$. Injective: if $2x_1 = 2x_2$ then $x_1 = x_2$.
        Surjective: any $y$ is hit by $x = y/2$. So bijective.
  \item $f(x) = x^2$ on $\R \to \R$. Not injective: $f(3) = f(-3) = 9$. Not surjective:
        $-4$ is never an output. Fails both.
  \item $f(x) = x^3$ on $\R \to \R$. Injective (cubing is strictly increasing) and
        surjective (every real has a real cube root). Bijective.
  \item $f(x) = e^x$ on $\R \to \R$. Injective but not surjective, as in the trap above.
\end{itemize}
\end{example}

\subsection{Why this is the spine of linear algebra}

Everything you will say about a matrix $\mat{A}$ being invertible is one of these three
words wearing a costume. Look at the map $T(\vec{x}) = \mat{A}\vec{x}$ for an $n\times n$
matrix.

\begin{itemize}
  \item $T$ is \textbf{injective} exactly when the only solution of $\mat{A}\vec{x} =
        \vec{0}$ is $\vec{x} = \vec{0}$ -- that is, when the null space is trivial, when
        the columns are independent. Two different inputs giving the same output would
        mean $\mat{A}(\vec{x}_1 - \vec{x}_2) = \vec{0}$ with a nonzero difference.
  \item $T$ is \textbf{surjective} exactly when the column space is all of $\R^n$ -- every
        target $\vec{b}$ is reachable, so $\mat{A}\vec{x} = \vec{b}$ always has a solution.
  \item $T$ is \textbf{bijective} exactly when $\mat{A}$ is invertible. Then, and only
        then, can you undo the machine.
\end{itemize}

So the entire ``invertibility'' story -- independent columns, trivial null space, full
column space, nonzero determinant -- is the demand that one linear machine be a bijection.
The abstract definitions you just met are the skeleton; the matrix theorems are the flesh.

\subsection{Inverses and composition}

When $f:A\to B$ is bijective, you can run it backwards. The \emph{inverse} $f^{-1}:B\to A$
sends each output back to the unique input that produced it, and it satisfies
\[
  f^{-1}(f(a)) = a \quad\text{and}\quad f(f^{-1}(b)) = b.
\]
Only bijections have inverses: you need surjectivity so that every $b$ has \emph{some}
preimage to return to, and injectivity so that the preimage is \emph{unique} (otherwise
$f^{-1}$ would not be a function). This is the abstract reason a non-square or
rank-deficient matrix has no inverse -- the underlying map fails one of the two halves.

\emph{Composition} chains two machines: if $f:A\to B$ and $g:B\to C$, then $g \circ f :
A \to C$ means ``do $f$, then $g$,'' so $(g\circ f)(a) = g(f(a))$. Read it right to left,
the way function notation forces (Figure~\ref{ch38:composition}). This is exactly how
matrix multiplication is born: applying $\mat{A}$ then $\mat{B}$ to a vector is the single
matrix $\mat{B}\mat{A}$, and the order matters because composition does not commute. The
notation $\mat{B}\mat{A}$ with $\mat{A}$ on the right is precisely $g\circ f$ with $f$ on
the right.

\begin{figure}[H]
\centering
\begin{tikzpicture}[scale=1.0,
   bag/.style={draw, ellipse, minimum width=1.5cm, minimum height=2.4cm, thick},
   dot/.style={circle,fill,inner sep=1.3pt}]
  \node[bag,sourceblue] (A) at (0,0) {};
  \node[bag,structurepurple] (B) at (3.2,0) {};
  \node[bag,targetgreen] (C) at (6.4,0) {};
  \node[below=2pt of A] {\small $A$};
  \node[below=2pt of B] {\small $B$};
  \node[below=2pt of C] {\small $C$};
  \node[dot,sourceblue] (a) at (0,0) {};
  \node[dot,structurepurple] (b) at (3.2,0) {};
  \node[dot,targetgreen] (c) at (6.4,0) {};
  \node[left=2pt of a] {\scriptsize $a$};
  \node[above=2pt of b] {\scriptsize $f(a)$};
  \node[right=2pt of c] {\scriptsize $g(f(a))$};
  \draw[-{Latex[length=2mm]}, sourceblue, thick] (a) to[bend left=22] node[above]{\scriptsize $f$} (b);
  \draw[-{Latex[length=2mm]}, targetgreen, thick] (b) to[bend left=22] node[above]{\scriptsize $g$} (c);
  \draw[-{Latex[length=2mm]}, quotientorange, thick, dashed] (a) to[bend right=28] node[below]{\scriptsize $g\circ f$} (c);
\end{tikzpicture}
\caption{Composition runs right to left: $g\circ f$ first applies $f$ (blue), landing in
$B$, then applies $g$ (green), landing in $C$. The dashed arrow is the single combined
machine $g\circ f$. Matrix multiplication $\mat{B}\mat{A}$ is this picture with $f=\mat{A}$
nearest the input.}
\label{ch38:composition}
\end{figure}

\begin{connectionbox}
The set of all bijections of a set to itself, under composition, is the prototype of a
\emph{group} -- composition is associative, identity is ``do nothing,'' and every
bijection has an inverse. When the bijections are \emph{linear}, this group is the general
linear group $GL_n(\R)$ of invertible matrices, the central object of the chapter on
groups and symmetry. So ``bijection plus composition'' is the seed from which all of group
theory grows, and invertible matrices are its most concrete flowering.
\end{connectionbox}

\begin{applicationbox}
In machine learning a \emph{neural network} is a long composition $f_L \circ \dots \circ
f_2 \circ f_1$ of simple maps, each a linear function followed by a wiggly nonlinearity.
Training adjusts the linear pieces; the chain rule of calculus -- itself a statement about
how derivatives compose -- propagates error backwards through the composition (this is
``backpropagation''). \emph{Normalizing flows} go further and insist every layer be a
\emph{bijection}, precisely so the network can be run backwards to generate data and so
that probabilities can be tracked exactly through the chain. Injectivity, surjectivity,
and composition are not museum pieces; they are design constraints in working systems.
\end{applicationbox}

\section{Fields: the numbers you are allowed to scale by}
\label{ch38:fields}

\subsection{The question hiding behind ``scalar''}

Every time you write $c\,\vec{v}$ you are scaling a vector by a number $c$. Which numbers
are legal? The main text answers ``real numbers, usually,'' and then, in the chapter on
complex matrices, ``complex numbers.'' But \emph{why those two}, and not the integers, or
just the positive numbers, or some other pool? Linear algebra
needs its scalars to support a specific list of guarantees -- and any number system
meeting that list is called a \emph{field}. $\R$ and $\C$ are simply the two fields you
already know and trust.

\begin{intuitionbox}
A field is a number system in which you can add, subtract, multiply, and
\emph{divide} by any nonzero element, with all the familiar rules of arithmetic holding.
The make-or-break ability is division: solving $c\,\vec{v} = \vec{w}$ for $\vec{v}$ means
multiplying by $1/c$, and the whole machinery of pivots, inverses, and bases depends on
being able to do that for any nonzero scalar. Fields are exactly the number systems where
``divide by the pivot'' never fails.
\end{intuitionbox}

\subsection{The definition, friendly first}

A field has two operations, $+$ and $\times$, and asks that each be as well-behaved as
possible. The cleanest way to say it: addition and multiplication both ``undo''
themselves.

\begin{formalbox}
A \textbf{field} is a set $\F$ with two operations $+$ and $\,\cdot\,$ such that for all
$a,b,c \in \F$:
\begin{itemize}
  \item \textbf{Addition} is commutative and associative, has an identity $0$, and every
        $a$ has a negative $-a$ with $a + (-a) = 0$.
  \item \textbf{Multiplication} is commutative and associative, has an identity $1 \ne 0$,
        and every \emph{nonzero} $a$ has an inverse $a^{-1}$ with $a \cdot a^{-1} = 1$.
  \item \textbf{Distributivity} links them: $a\cdot(b+c) = a\cdot b + a\cdot c$.
\end{itemize}
\end{formalbox}

The single line with real teeth is ``every nonzero $a$ has an inverse $a^{-1}$.'' That is
division. Everything else on the list you have believed since grade school. Compare with
the vector-space axioms in the main text: they have the same shape, and that is no
accident -- a field is, among other things, a tiny one-dimensional vector space over
itself, and the field axioms are exactly the guarantees the scalars must provide so the
vector-space axioms can do their work.

\begin{historybox}
In May 1832, on the eve of his fatal duel at age twenty, the brilliant French mathematician \'Evariste Galois feverishly scribbled his mathematical testament in a letter to his friend Auguste Chevalier. Galois was solving a mystery that had baffled algebraists for centuries: why can quadratic, cubic, and quartic polynomial equations be solved by formulas with radicals ($\sqrt{\cdot}$), while the general quintic equation cannot? Galois saw that roots do not live in isolation. They generate an algebraic number field, and the symmetries of these numbers under permutations form a group. Galois had submitted earlier memoirs to the French Academy of Sciences, but Augustin-Louis Cauchy lost them, Joseph Fourier died with one manuscript on his desk, and Sim\'eon Denis Poisson rejected another as ``incomprehensible.'' Only decades after his death was Galois's vision vindicated.

In 1910, German mathematician Ernst Steinitz published \emph{Algebraische Theorie der K\"orper} (The Algebraic Theory of Fields), establishing the modern axiomatic foundation for fields. Steinitz showed that fields naturally split by their characteristic: infinite fields like the rational numbers $\mathbb{Q}$ and the real numbers $\R$, and finite fields $\mathbb{F}_p = \Z/p\Z$ built from clock arithmetic modulo a prime $p$.

Why does this matter for linear algebra? Because virtually every algebraic theorem in this book---Gaussian elimination, rank-nullity, matrix inversion, determinants, and subspace dimensions---relies only on the field axioms. It does not matter whether the scalars are real measurements or finite integers modulo 2. While continuous physics and geometry require $\R$ and $\C$, modern computer science, error-correcting codes, and cryptography (such as the AES block cipher and elliptic curves) operate entirely on vector spaces over finite fields $\mathbb{F}_{2^8}$ or $\mathbb{F}_p$.
\end{historybox}

\subsection{Who passes, who fails}

Run the systems we tabulated earlier past the field test, watching the division line.

\begin{center}
\begin{tabular}{lcl}
\toprule
System & Field? & Reason \\
\midrule
$\R$ (reals) & yes & add, subtract, multiply, divide by nonzero -- all fine \\
$\C$ (complex) & yes & same, and even better: every polynomial has a root \\
$\Q$ (rationals) & yes & fractions of fractions are fractions \\
$\Z$ (integers) & \textbf{no} & $1/2 \notin \Z$: cannot divide \\
$\N$ (naturals) & \textbf{no} & no negatives, no division \\
\bottomrule
\end{tabular}
\end{center}

This is the precise reason linear algebra ``over the integers'' is harder and stranger:
$\Z$ is not a field, so you cannot freely divide by a pivot, and Gaussian elimination
stalls. The instant you allow yourself fractions -- step up to $\Q$ or $\R$ -- elimination
flows, every nonzero pivot is invertible, and the clean theory you know returns. When the
main text says ``let the scalars be real,'' it is really saying ``let the scalars be a
field, and $\R$ is the one we picture.''

\begin{example}[A finite field: arithmetic on a clock]\label{ch38:f2}
Fields need not be infinite. Take $\F_2 = \{0, 1\}$ with addition and multiplication done
``modulo $2$'' -- keep only the remainder after dividing by $2$:
\[
  0+0=0,\quad 0+1=1,\quad 1+1=0,\qquad
  0\cdot 0 = 0,\quad 0\cdot 1 = 0,\quad 1\cdot 1 = 1.
\]
Check the teeth: the only nonzero element is $1$, and $1\cdot 1 = 1$, so $1$ is its own
inverse. Division works. $\F_2$ is a genuine field with just two elements -- it is the
arithmetic of parity, of a single bit, of XOR. Linear algebra over $\F_2$ is the
mathematics of error-correcting codes and cryptography, and every theorem about bases and
rank you prove ``over a field'' holds there verbatim.
\end{example}

\begin{aestheticbox}
Here is what the field abstraction buys you. You prove the theory of bases,
dimension, rank, and linear maps \emph{once}, assuming only ``the scalars form a field.''
That single proof then pays out over $\R$ (geometry and data), over $\C$ (signals,
quantum mechanics, the spectral theorem), over $\Q$ (exact symbolic computation), and over
$\F_2$ and its finite cousins (coding theory, cryptography). One abstraction, four wildly
different worlds, zero extra work. That economy -- prove it where it is true in general,
harvest it everywhere -- is the deepest habit of mind linear algebra teaches.
\end{aestheticbox}

\subsection{Why $\R$ and why $\C$}

Among all fields, two earn their place as the default scalars for very human reasons.

$\R$ is the field of \emph{measurement}. Lengths, angles, masses, probabilities,
coordinates of a point -- the physical world is reported in real numbers, and the real
line is \emph{complete} (no gaps), which is exactly what makes limits, continuity, and
hence calculus possible. When linear algebra meets the world, it meets it over $\R$.

$\C$ is the field of \emph{completion}. The reals have a flaw: the innocent equation
$x^2 + 1 = 0$ has no real solution. Adjoining a single new number $i$ with $i^2 = -1$
repairs this flaw \emph{totally} -- over $\C$ every polynomial of degree $n$ has exactly
$n$ roots (the Fundamental Theorem of Algebra). For linear algebra this is not a
curiosity; it is the reason every square matrix has a full set of eigenvalues over $\C$,
why rotations have ``imaginary'' eigenvalues, and why the spectral story is cleanest in
the complex setting. $\R$ is where the data lives; $\C$ is where the algebra is whole.

\begin{connectionbox}
The two reasons above point at two different branches of mathematics. The completeness of
$\R$ is an \emph{analysis} fact -- it is what lets you take limits, and it underwrites
every convergence argument in numerical linear algebra and in the calculus that powers
optimization. The algebraic closure of $\C$ is an \emph{algebra} fact -- it is what lets
every characteristic polynomial factor completely, and it is the bridge to the theory of
eigenvalues, Jordan form, and beyond. Choosing $\R$ or $\C$ as your scalars quietly
decides which of these two great traditions you are leaning on.
\end{connectionbox}

\section{Just enough proof literacy}

The main chapters prove things, gently but for real. You do not need to be a proof athlete
to follow them, but a few patterns recur so often that recognizing them turns proofs from
intimidating to almost predictable. This section is a field guide to those patterns.

\subsection{Reading the logical connectives}

Three little phrases do most of the work in any proof.

\begin{itemize}
  \item \textbf{Implication} $P \implies Q$: ``if $P$, then $Q$.'' To prove it, you assume
        $P$ is true and argue your way to $Q$. You never have to worry about the case
        where $P$ is false -- the claim says nothing there.
  \item \textbf{Equivalence} $P \iff Q$: ``$P$ if and only if $Q$,'' often abbreviated
        ``iff.'' This is two implications bundled: $P \implies Q$ \emph{and} $Q \implies
        P$. To prove an iff you must prove \emph{both} directions, and forgetting one is
        the most common gap in a beginner's proof.
  \item \textbf{Quantifiers} ``for all'' ($\forall$) and ``there exists'' ($\exists$).
        ``For all'' is a promise about every case; to \emph{disprove} it you need just one
        counterexample. ``There exists'' is a claim that at least one case works; to prove
        it you exhibit one witness, and to disprove it you must rule out all cases.
\end{itemize}

\begin{trapbox}
$P \implies Q$ does \emph{not} give you $Q \implies P$. ``If $\vec{x}$ is the zero vector,
then $\mat{A}\vec{x} = \vec{0}$'' is always true; its reverse ``if $\mat{A}\vec{x} =
\vec{0}$ then $\vec{x} = \vec{0}$'' is true only for special matrices (the injective
ones). Confusing an implication with its converse is the single most frequent logical slip
in linear algebra. When a theorem is an ``iff,'' respect that it is making two separate
promises.
\end{trapbox}

\subsection{The handful of proof shapes}

Almost every proof in this book is one of four moves.

\textbf{Direct proof.} Assume the hypothesis, walk forward through definitions and
algebra, arrive at the conclusion. Most ``closure'' arguments are direct: to show a set is
a subspace you take two members, add them, and check the sum is still a member -- one
short computation.

\textbf{Proof by contradiction.} Suppose the claim is \emph{false}, then derive something
impossible (like $0 = 1$, or a vector that is both zero and nonzero). The impossibility
means your supposition was wrong, so the claim is true. The classic example is showing a
set of vectors is independent: assume a nontrivial combination equals zero, then squeeze
out a contradiction.

\textbf{Proof by contrapositive.} To prove $P \implies Q$, prove instead the logically
identical $\lnot Q \implies \lnot P$ (``not $Q$ implies not $P$''). These two statements
are \emph{always} equivalent, so you may prove whichever is easier. Often the
contrapositive is the friendlier road.

\textbf{Proof by induction.} To prove a statement for all natural numbers $n$: show it for
the base case $n = 0$ (or $1$), then show that its truth at $n$ forces its truth at $n+1$.
Like dominoes -- knock over the first, guarantee each knocks the next, and they all fall.
Dimension arguments and ``every matrix can be reduced in finitely many steps'' arguments
often run on induction.

\begin{example}[A direct proof you can already follow]\label{ch38:directproof}
\emph{Claim.} If $f:A\to B$ and $g:B\to C$ are both injective, then $g\circ f : A \to C$
is injective.

\emph{Proof.} We use the proof-friendly form of injectivity: ``equal outputs force equal
inputs.'' Suppose $(g\circ f)(a_1) = (g\circ f)(a_2)$, i.e.\ $g(f(a_1)) = g(f(a_2))$. Since
$g$ is injective, equal outputs of $g$ force equal inputs, so $f(a_1) = f(a_2)$. Since $f$
is injective, this in turn forces $a_1 = a_2$. We assumed equal outputs of $g\circ f$ and
deduced equal inputs, so $g \circ f$ is injective. $\qquad\blacksquare$
\end{example}

Read that proof twice. Nothing in it is clever; it is just the definitions, applied in
order, with each ``since'' pointing back to exactly one hypothesis. The vast majority of
proofs in this book have that same unhurried texture. When one feels hard, slow down and
ask of each sentence: which definition or hypothesis is this using? Almost always the
answer is sitting one line up.

\begin{walkbox}
On your next walk, take the claim ``if $g \circ f$ is injective, then $f$ is injective''
(note: only $f$, not $g$). Is it true? Try to either prove it directly or build a tiny
counterexample with three- or four-element sets and arrows on paper. Then ask the mirror
question about surjectivity: if $g \circ f$ is surjective, which of $f$, $g$ must be
surjective? Getting these straight in your head, away from a screen, cements the whole
section.
\end{walkbox}

\begin{strangbox}
Strang's lectures move fast through this background -- he assumes you can read set-builder
notation and knows you will pick up ``onto'' and ``one-to-one'' from context.
\emph{Before} Strang, this appendix is the gentle on-ramp: get comfortable here and his
definitions will feel like old friends rather than new hurdles. \emph{After} Strang,
return for the field perspective -- once you have seen eigenvalues fail to be real and then
reappear over $\C$, the ``why $\R$ and why $\C$'' section explains the pattern you noticed
but he did not pause to name.
\end{strangbox}

\begin{focusbox}
The whole appendix in one breath. A \textbf{set} answers ``in or out.'' A
\textbf{function} is a deterministic machine; \textbf{injective} means no output is reused,
\textbf{surjective} means no output is missed, and \textbf{bijective} (both) means it can
be run backwards -- which for a matrix is exactly invertibility. A \textbf{field} is a
number system where you can divide by anything nonzero; $\R$ and $\C$ are the two we use
because one measures the world and the other completes the algebra. And a \textbf{proof}
is almost always just the definitions applied in order. Carry these and the main text will
rarely surprise you.
\end{focusbox}

\subsection*{What We Can Now Do (Capabilities Unlocked)}
\begin{itemize}
  \item \textbf{Axiomatic Substrate for Vector Spaces:} We formalize the precise mathematical requirements for vectors, sets, and scalars, understanding why vector spaces can be defined over any arbitrary field $\F$.
  \item \textbf{Function-Matrix Correspondence:} We connect the abstract set-theoretic properties of functions (injectivity, surjectivity, bijectivity) directly to fundamental linear algebra properties (trivial nullspace, full column rank, and two-sided invertibility).
  \item \textbf{Finite Field Linear Algebra:} By abstracting away from real Euclidean geometry, we can apply matrix multiplication, Gaussian elimination, and rank-nullity to finite fields $\mathbb{F}_p$ and Galois fields $\mathbb{F}_{2^k}$, powering computer arithmetic, hash functions, and modern cryptography.
  \item \textbf{Proof Mechanics:} We turn abstract definitions into mechanical proof recipes: unpacking definitions, verifying hypotheses, and constructing direct or contrapositive arguments.
\end{itemize}

\subsection*{Honest Downsides and Limitations}
\begin{itemize}
  \item \textbf{Loss of Metric Geometry over Finite Fields:} Over finite fields $\mathbb{F}_p$, there is no notion of order ($x > 0$ is meaningless), and lengths or dot products behave strangely. For example, in $\mathbb{F}_2^2$, the vector $\colvec{1\\1}$ has dot product $1\cdot 1 + 1\cdot 1 = 1 + 1 = 0 \pmod 2$, meaning a nonzero vector is orthogonal to itself! Concepts like angles, positive definiteness, and SVD rely heavily on $\R$ and $\C$.
  \item \textbf{Russell's Paradox in Naive Set Theory:} Naive set comprehension ($\{x \mid P(x)\}$) allows self-referential paradoxes like the set of all sets that do not contain themselves ($R = \{x \mid x \notin x\}$), requiring the heavy machinery of Zermelo--Fraenkel axiomatic set theory to resolve.
  \item \textbf{Lack of Algorithmic Complexity:} Set theory and abstract field theory describe algebraic structures, but do not provide efficient computational algorithms for calculating large determinants or factoring polynomials.
\end{itemize}

\subsection*{Where We Go Next in the Book}
\begin{itemize}
  \item \textbf{Appendix B (Chapter 39: Sums, Complex Numbers, and Calculus):} With sets and fields secured, we turn to the concrete computational machinery: index summation techniques ($\Sigma$), complex numbers and polar Euler forms ($e^{i\theta}$), and multivariable matrix calculus.
\end{itemize}

\section{Exercises}

\begin{exercise}[Reading builder notation]
Translate each set into plain English and list its elements (or describe them) explicitly.
\[
  \text{(a) } \set{n \in \Z : -3 \le n \le 2}, \qquad
  \text{(b) } \set{x \in \R : x^2 = 4}, \qquad
  \text{(c) } \set{(x,y)\in\R^2 : x + y = 1}.
\]
For (c), what geometric object is it?
\end{exercise}

\begin{exercise}[Set combinations]
Let $A = \{1,2,3,4,5\}$ and $B = \{2,4,6\}$. Compute $A \cup B$, $A \cap B$, $A \setminus
B$, and $B \setminus A$. Then verify on this example that $A \setminus B$ and $B \setminus
A$ are generally different sets.
\end{exercise}

\begin{exercise}[Classifying functions]
For each function decide whether it is injective, surjective, both, or neither, and justify
in one sentence. Take domain and codomain $\R \to \R$ unless stated.
\[
  \text{(a) } f(x) = 3x - 1, \quad
  \text{(b) } f(x) = x^2 + 1, \quad
  \text{(c) } f(x) = x^3 - x, \quad
  \text{(d) } f:\R\to[0,\infty),\ f(x)=x^2.
\]
\end{exercise}

\begin{exercise}[Codomain matters]
The formula $f(x) = \tan(x)$ on the open interval $(-\pi/2, \pi/2)$ can be viewed as a map
into $\R$. Argue that it is bijective onto $\R$, and explain in words why shrinking the
domain (say to $(0,\pi/2)$) would ruin surjectivity. What does this say in general about
the two loopholes a function has?
\end{exercise}

\begin{exercise}[Inverses need bijections]
Explain, using the definitions of injective and surjective, exactly which property fails
when you try to invert $f(x) = x^2$ on $\R\to\R$. Then show that restricting both domain
and codomain to $[0,\infty)$ makes it bijective, and find the inverse formula.
\end{exercise}

\begin{exercise}[Is it a field?]
Decide whether each is a field, and if not, name the axiom(s) that fail.
\[
  \text{(a) } \Q, \qquad
  \text{(b) the even integers } 2\Z, \qquad
  \text{(c) } \set{a + b\sqrt 2 : a,b \in \Q}.
\]
For (c) you may take on faith that sums and products stay in the set; focus on whether
nonzero elements have inverses. (Hint for (c): rationalize $\dfrac{1}{a + b\sqrt 2}$.)
\end{exercise}

\begin{exercise}[The two-element field in action]
In $\F_2 = \{0,1\}$ with arithmetic modulo $2$, solve the linear system
\[
  x + y = 1, \qquad y + z = 0, \qquad x + z = 1
\]
for $x,y,z \in \F_2$. (Remember $1+1 = 0$.) How many solutions are there? Contrast this
with how the same-shaped system would behave over $\R$.
\end{exercise}

\begin{exercise}[Implication versus converse]
For a matrix $\mat{A}$, consider $P$: ``the columns of $\mat{A}$ are linearly independent''
and $Q$: ``$\mat{A}\vec{x} = \vec{b}$ has a solution for every $\vec{b}$.'' For a general
(non-square) matrix, can $P$ hold while $Q$ fails? Can $Q$ hold while $P$ fails? For a
\emph{square} matrix, explain why they become equivalent. (You may quote results from the
main text; the point is to practice separating an implication from its converse.)
\end{exercise}

\begin{exercise}[A composition proof]
Prove, in the unhurried style of Example~\ref{ch38:directproof}: if $f:A\to B$ and
$g:B\to C$ are both surjective, then $g\circ f$ is surjective. (Start from an arbitrary
$c \in C$ and walk backwards through the two surjectivity assumptions.)
\end{exercise}

\begin{exercise}[Counterexample hunting]
Decide whether each statement is true; if false, give the smallest counterexample you can
(small finite sets and arrows on paper are ideal).
\begin{itemize}
  \item[(a)] If $g\circ f$ is injective, then $g$ is injective.
  \item[(b)] If $g\circ f$ is injective, then $f$ is injective.
  \item[(c)] If $g\circ f$ is surjective, then $g$ is surjective.
\end{itemize}
\end{exercise}

\begin{exercise}[Connecting to the matrix story]
Let $T:\R^3 \to \R^2$ be $T(\vec{x}) = \mat{A}\vec{x}$ for some $2\times 3$ matrix
$\mat{A}$. Using the function definitions only, explain why $T$ \emph{cannot} be injective
no matter what $\mat{A}$ is, but \emph{can} be surjective. Relate your reasoning to the
sizes $3$ (inputs) and $2$ (outputs), and to the idea of a null space.
\end{exercise}
